\section{Tổng quan về bài toán và phạm vi đề tài}

\subsection{Bối cảnh và mục tiêu}

Nhiếp ảnh phim đang quay trở lại như một trào lưu sáng tác nghiêm túc, kéo theo nhu cầu tổ chức các cuộc thi dành riêng cho ảnh chụp trên phim. Khác với cuộc thi ảnh số, ban tổ chức ở đây phải xác minh được rằng tác phẩm thực sự được chụp bằng phim thật, không phải ảnh số hay ảnh do trí tuệ nhân tạo sinh ra.

Ảnh số mang sẵn siêu dữ liệu EXIF ghi lại máy ảnh, ống kính và thông số phơi sáng. Hầu hết máy ảnh phim không ghi lại gì, nên toàn bộ thông tin kỹ thuật đều do thí sinh tự khai. Bài toán vì vậy không chỉ là lưu trữ mà còn là kiểm chứng, và đó là lý do hệ thống cần một cơ sở dữ liệu được thiết kế chặt chẽ thay vì vài bảng dựng theo cảm tính.

\subsection{Phạm vi báo cáo bốn chương}

Báo cáo đi theo đúng chu kỳ sống của một cơ sở dữ liệu, mỗi chương nhận đầu ra của chương trước làm đầu vào:


\begin{itemize}
    \item Chương 1 khảo sát nghiệp vụ và phát biểu tập quy tắc nghiệp vụ, kết thúc bằng bảng ánh xạ yêu cầu sang đối tượng dữ liệu.
    \item Chương 2 xây dựng mô hình quan niệm: thực thể, thuộc tính, mối quan hệ, bản số và sơ đồ ERD.
    \item Chương 3 chuyển sang mô hình logic, phân tích phụ thuộc hàm, xác định khóa ứng viên, chuẩn hóa và lập từ điển dữ liệu.
    \item Chương 4 thiết kế vật lý: ánh xạ kiểu dữ liệu, đánh chỉ mục, viết mã DDL, phân tích tối ưu truy vấn và xây dựng các truy vấn mẫu.
\end{itemize}


